\chapter{Painting: Open Causal Spaces and Benign Underdetermination}
\label{ch:painting}

\begin{quotation}
\noindent\textit{Synopsis.} This chapter examines painting as the canonical
open hidden-state domain. Paintings routinely omit vast amounts of
information while remaining meaningful and useful; spatial compression,
implied motion, and inferred causal history are analyzed as cases where a
viewer reconstructs a rich world from sparse cues. Unlike chess, no finite
augmentation can enumerate every causal history consistent with a painted
configuration. The chapter argues that ambiguity here is not merely
tolerated but often serves the intended task directly, making painting the
clearest illustration in the monograph of the difference between
corruption and benign underdetermination.
\end{quotation}

\section{Spatial Compression}

A small number of brushstrokes suggesting a cluster of leaves, and a
larger arrangement of such clusters suggesting a forest, exploit a
viewer's priors under \cref{def:benign-loss}: the omitted detail is
redundant given what the visual system already supplies, not required
for the viewer's task of recognizing foliage.

\begin{remark}[Priors as Repair Cost Reduction]
\label{rem:priors-reduce-cost}
It is worth stating this connection to Chapter~\ref{ch:repair-economics}
explicitly rather than leaving spatial compression and repair economics
as two separate vocabularies that happen to apply to the same example. A
painter can afford to discard a leaf-level distinction $d$ at negligible
$\repcost{d}$ precisely because the viewer's own perceptual system
performs, for free and at the moment of viewing, whatever reconstruction
would otherwise have been needed to supply $d$. The painter's low repair
investment is not recklessness; it is an accurate assessment that the
\emph{effective} cost of maintaining $d$ for this consumer's task is
already near zero, since the consumer's prior does the work the painter
would otherwise have had to fund.
\end{remark}

\section{Temporal Compression and the Bent Branch}

A branch bent as if by wind is a static configuration that is
low-probability under a viewer's prior model of objects at rest, forcing
an inference of cause. Informally,
\[
P(\text{cause} \mid \text{state}) \propto P(\text{state} \mid
\text{cause})\, P(\text{cause}),
\]
used here descriptively, to fix intuition, rather than as a proved result
of this monograph.

\begin{proposition}[Open Causal Space]
\label{prop:open-causal-space}
No finite augmentation of a painting's visible configuration determines
a unique causal history for that configuration.
\end{proposition}

\begin{proofsketch}
For any finite augmentation set $H$ added to the image, multiple distinct
causal histories -- wind, a bird's departure, a hand's release,
accumulated then-melted snow -- remain consistent with the image together
with $H$. The causal space is therefore open in the sense of
\cref{def:closed-open}, and no finite $H$ closes it.
\end{proofsketch}

\begin{proposition}[Benign Ambiguity]
\label{prop:benign-ambiguity}
Multiple admissible causal histories for a painted configuration need not
reduce success at the task of evocative reconstruction.
\end{proposition}

\begin{proofsketch}
The task of evocative reconstruction (\cref{def:sufficiency} applied with
$\tau = $ ``evoke a plausible dynamic scene'') does not require a unique
history; any history in $\admis{\tau}{r}$ satisfies $\tau$ equally well.
Ambiguity here is benign in the sense of \cref{def:benign-loss} precisely
because the task does not discriminate among the remaining admissible
continuations. By \cref{prop:benign-preserves-fidelity}, this ambiguity
is also compatible with $\fidelity{r}{\tau} = 1$ for this task, even
though the same rendering would have far lower fidelity for a task
demanding the branch's one true cause.
\end{proofsketch}

\section{Why Painting Is the Canonical Open Domain}

\cref{prop:open-causal-space} is the direct converse witness to
\cref{prop:sufficiency-closed-only}: painting demonstrates that
sufficiency by augmentation is unavailable precisely where the hidden
state is open, exactly as chess (\cref{ch:chess}) demonstrates its
availability where the hidden state is closed.

\section{A Preview: Closedness and Declaration Are Independent}

It is tempting, having seen chess (closed, and its hidden state
straightforwardly enumerable) and painting (open, and its causal space
irreducibly ambiguous), to treat closedness and transparency as the same
property. Chapter~\ref{ch:maps} shows they are not: a map depicts a
domain (terrain) that is, in its full physical complexity, at least as
rich as anything a painting depicts, yet a map's discard set is closed
and declared by its scale bar. The relevant two-by-two structure is worth
stating here in preview, since painting and chess between them supply
only two of its four cells.

\begin{center}
\begin{tabular}{lll}
\toprule
& Declared discard & Undeclared discard \\
\midrule
Closed hidden state & Maps (\cref{ch:maps}) & Chess without $H$ \\
Open hidden state & A destroyed-records archive & Painting (this chapter) \\
\bottomrule
\end{tabular}
\end{center}

The upper-right cell is a chess position reported without its hidden
state -- closed, but undeclared, and hence merely incomplete rather than
irreducibly ambiguous, since $H$ could in principle be supplied. The
lower-left cell is populated by a case not otherwise treated in this
monograph: an archive that explicitly records that certain materials were
destroyed, without being able to enumerate what they contained -- open,
since the destroyed content's specifics are gone rather than merely
omitted from a known list, but declared, since the fact and approximate
scope of the loss is a matter of record. Painting occupies the remaining
cell: open, and undeclared in the specific sense that a painting does not
carry an explicit ledger of which causal histories it fails to
distinguish, relying instead on the viewer's tolerance for ambiguity
(\cref{prop:benign-ambiguity}) rather than on disclosure.
Chapter~\ref{ch:maps} develops the closed-and-declared cell in full.
